\documentclass[11pt,a4paper]{article}
% ==============================================================
%  Préambule commun aux fiches -- Advanced Probabilities (M1)
%  Usage, dans chaque fiche :
%     \documentclass[11pt,a4paper]{article}
%     % ==============================================================
%  Préambule commun aux fiches -- Advanced Probabilities (M1)
%  Usage, dans chaque fiche :
%     \documentclass[11pt,a4paper]{article}
%     % ==============================================================
%  Préambule commun aux fiches -- Advanced Probabilities (M1)
%  Usage, dans chaque fiche :
%     \documentclass[11pt,a4paper]{article}
%     \input{../preambule.tex}
%     \fiche{C5}{Vecteurs gaussiens}{Gaussian vectors}
% ==============================================================
\usepackage[T1]{fontenc}
\usepackage[utf8]{inputenc}
\usepackage[french]{babel}
\usepackage{lmodern}
\usepackage[a4paper,margin=1.9cm,top=1.7cm,bottom=1.9cm]{geometry}
\usepackage{amsmath,amssymb,amsthm,mathtools}
\usepackage[inline]{enumitem}
\usepackage{booktabs,array,tabularx}
\usepackage{xcolor}
\usepackage[most]{tcolorbox}
\usepackage{fancyhdr}
\usepackage{hyperref}
\hypersetup{colorlinks=true,linkcolor=blue!50!black,urlcolor=blue!50!black}

\setlist{itemsep=1pt,topsep=3pt,parsep=0pt}
\setlength{\parindent}{0pt}
\setlength{\parskip}{3pt}
\emergencystretch=3em

% ---------- Identification de la fiche -------------------------
\newcommand{\ficheid}{}
\newcommand{\fichetitre}{}
\newcommand{\fiche}[3]{%
  \renewcommand{\ficheid}{#1}\renewcommand{\fichetitre}{#2}%
  \pagestyle{fancy}\fancyhf{}%
  \fancyhead[L]{\small\textsf{Advanced Probabilities -- Fiche #1}}%
  \fancyhead[R]{\small\textsf{#2}}%
  \fancyfoot[C]{\small\thepage}%
  \begin{center}
    {\LARGE\bfseries Fiche #1 -- #2}\par\smallskip
    {\large\itshape (#3)}\par\smallskip
    {\small Cours de V.\,M.~Lambert, M1 -- semestre 1}
  \end{center}\vspace{-2pt}\hrule\vspace{6pt}}

% ---------- Environnements numérotés (un seul compteur) --------
\theoremstyle{definition}
\newtheorem{definition}{Définition}
\renewcommand{\thedefinition}{\ficheid.\arabic{definition}}
\newtheorem{theoreme}[definition]{Théorème}
\newtheorem{proposition}[definition]{Proposition}
\newtheorem{lemme}[definition]{Lemme}
\newtheorem{corollaire}[definition]{Corollaire}
\newtheorem{exemple}[definition]{Exemple}
\theoremstyle{remark}
\newtheorem{remarque}[definition]{Remarque}

% Résultat hors cours (standard, signalé comme tel)
\theoremstyle{definition}
\newtheorem{horscours}[definition]{Hors cours}

% Idée de preuve
\newenvironment{idee}{\par\small\textit{Idée de preuve.}~}{\par}

% ---------- Encadrés -----------------------------------------
\tcbset{fichebox/.style={enhanced,breakable,boxrule=0.6pt,arc=2pt,
  left=5pt,right=5pt,top=3pt,bottom=3pt,fonttitle=\bfseries}}
\newtcolorbox{piege}[1][]{fichebox,colback=red!4,colframe=red!60!black,
  title={Pièges classiques},#1}
\newtcolorbox{liens}[1][]{fichebox,colback=blue!3,colframe=blue!50!black,
  title={Liens entre notions},#1}
\newtcolorbox{methode}[1][Méthode]{fichebox,colback=green!4,colframe=green!40!black,title={#1}}
\newtcolorbox{alerte}[1][]{fichebox,colback=orange!6,colframe=orange!70!black,
  title={Attention à l'énoncé},#1}

% Exercice de TD (fiches de rappels)
\newcounter{exo}
\newenvironment{exo}[1][]{\refstepcounter{exo}\par\medskip
  \noindent{\large\bfseries Exercice~\theexo}\ifx\relax#1\relax\else~{\large\bfseries -- #1}\fi\par\nopagebreak\smallskip}{\par}
\newcommand{\outils}{\textbf{Outils du cours.}~}
\newcommand{\meth}{\textbf{Méthode / astuce.}~}
\newcommand{\indic}{\textbf{Indication.}~}
\newcommand{\rep}{\textbf{Réponse.}~}

% ---------- Notations du cours ---------------------------------
\newcommand{\R}{\mathbb{R}}
\newcommand{\N}{\mathbb{N}}
\newcommand{\Q}{\mathbb{Q}}
\newcommand{\C}{\mathbb{C}}
\newcommand{\PP}{\mathbb{P}}
\newcommand{\QQ}{\mathbb{Q}}
\newcommand{\E}{\mathbb{E}}
\newcommand{\F}{\mathcal{F}}
\newcommand{\G}{\mathcal{G}}
\newcommand{\Hc}{\mathcal{H}}
\newcommand{\Tc}{\mathcal{T}}
\newcommand{\Bc}{\mathcal{B}}
\newcommand{\Cc}{\mathcal{C}}
\newcommand{\Ec}{\mathcal{E}}
\newcommand{\Pc}{\mathcal{P}}
\newcommand{\Nc}{\mathcal{N}}
\newcommand{\LN}{\mathcal{LN}}
\newcommand{\Expo}{\mathcal{E}}
\newcommand{\Poi}{\mathcal{P}}
\newcommand{\Ff}{\mathbb{F}}
\newcommand{\1}{\mathbf{1}}
\newcommand{\dd}{\,\mathrm{d}}
\DeclareMathOperator{\Var}{\mathbb{V}}
\DeclareMathOperator{\Cov}{Cov}
\DeclareMathOperator{\card}{card}
\DeclareMathOperator{\diag}{diag}
\newcommand{\indep}{\perp\!\!\!\perp}
\newcommand{\ps}{\ \PP\text{-p.s.}}
\newcommand{\tendloi}{\xrightarrow[n\to+\infty]{\text{loi}}}
\newcommand{\tendps}{\xrightarrow[n\to+\infty]{\text{p.s.}}}
\newcommand{\tendproba}{\xrightarrow[n\to+\infty]{\PP}}
\newcommand{\tendLp}{\xrightarrow[n\to+\infty]{L^p}}
\newcommand{\tend}{\xrightarrow[n\to+\infty]{}}

%     \fiche{C5}{Vecteurs gaussiens}{Gaussian vectors}
% ==============================================================
\usepackage[T1]{fontenc}
\usepackage[utf8]{inputenc}
\usepackage[french]{babel}
\usepackage{lmodern}
\usepackage[a4paper,margin=1.9cm,top=1.7cm,bottom=1.9cm]{geometry}
\usepackage{amsmath,amssymb,amsthm,mathtools}
\usepackage[inline]{enumitem}
\usepackage{booktabs,array,tabularx}
\usepackage{xcolor}
\usepackage[most]{tcolorbox}
\usepackage{fancyhdr}
\usepackage{hyperref}
\hypersetup{colorlinks=true,linkcolor=blue!50!black,urlcolor=blue!50!black}

\setlist{itemsep=1pt,topsep=3pt,parsep=0pt}
\setlength{\parindent}{0pt}
\setlength{\parskip}{3pt}
\emergencystretch=3em

% ---------- Identification de la fiche -------------------------
\newcommand{\ficheid}{}
\newcommand{\fichetitre}{}
\newcommand{\fiche}[3]{%
  \renewcommand{\ficheid}{#1}\renewcommand{\fichetitre}{#2}%
  \pagestyle{fancy}\fancyhf{}%
  \fancyhead[L]{\small\textsf{Advanced Probabilities -- Fiche #1}}%
  \fancyhead[R]{\small\textsf{#2}}%
  \fancyfoot[C]{\small\thepage}%
  \begin{center}
    {\LARGE\bfseries Fiche #1 -- #2}\par\smallskip
    {\large\itshape (#3)}\par\smallskip
    {\small Cours de V.\,M.~Lambert, M1 -- semestre 1}
  \end{center}\vspace{-2pt}\hrule\vspace{6pt}}

% ---------- Environnements numérotés (un seul compteur) --------
\theoremstyle{definition}
\newtheorem{definition}{Définition}
\renewcommand{\thedefinition}{\ficheid.\arabic{definition}}
\newtheorem{theoreme}[definition]{Théorème}
\newtheorem{proposition}[definition]{Proposition}
\newtheorem{lemme}[definition]{Lemme}
\newtheorem{corollaire}[definition]{Corollaire}
\newtheorem{exemple}[definition]{Exemple}
\theoremstyle{remark}
\newtheorem{remarque}[definition]{Remarque}

% Résultat hors cours (standard, signalé comme tel)
\theoremstyle{definition}
\newtheorem{horscours}[definition]{Hors cours}

% Idée de preuve
\newenvironment{idee}{\par\small\textit{Idée de preuve.}~}{\par}

% ---------- Encadrés -----------------------------------------
\tcbset{fichebox/.style={enhanced,breakable,boxrule=0.6pt,arc=2pt,
  left=5pt,right=5pt,top=3pt,bottom=3pt,fonttitle=\bfseries}}
\newtcolorbox{piege}[1][]{fichebox,colback=red!4,colframe=red!60!black,
  title={Pièges classiques},#1}
\newtcolorbox{liens}[1][]{fichebox,colback=blue!3,colframe=blue!50!black,
  title={Liens entre notions},#1}
\newtcolorbox{methode}[1][Méthode]{fichebox,colback=green!4,colframe=green!40!black,title={#1}}
\newtcolorbox{alerte}[1][]{fichebox,colback=orange!6,colframe=orange!70!black,
  title={Attention à l'énoncé},#1}

% Exercice de TD (fiches de rappels)
\newcounter{exo}
\newenvironment{exo}[1][]{\refstepcounter{exo}\par\medskip
  \noindent{\large\bfseries Exercice~\theexo}\ifx\relax#1\relax\else~{\large\bfseries -- #1}\fi\par\nopagebreak\smallskip}{\par}
\newcommand{\outils}{\textbf{Outils du cours.}~}
\newcommand{\meth}{\textbf{Méthode / astuce.}~}
\newcommand{\indic}{\textbf{Indication.}~}
\newcommand{\rep}{\textbf{Réponse.}~}

% ---------- Notations du cours ---------------------------------
\newcommand{\R}{\mathbb{R}}
\newcommand{\N}{\mathbb{N}}
\newcommand{\Q}{\mathbb{Q}}
\newcommand{\C}{\mathbb{C}}
\newcommand{\PP}{\mathbb{P}}
\newcommand{\QQ}{\mathbb{Q}}
\newcommand{\E}{\mathbb{E}}
\newcommand{\F}{\mathcal{F}}
\newcommand{\G}{\mathcal{G}}
\newcommand{\Hc}{\mathcal{H}}
\newcommand{\Tc}{\mathcal{T}}
\newcommand{\Bc}{\mathcal{B}}
\newcommand{\Cc}{\mathcal{C}}
\newcommand{\Ec}{\mathcal{E}}
\newcommand{\Pc}{\mathcal{P}}
\newcommand{\Nc}{\mathcal{N}}
\newcommand{\LN}{\mathcal{LN}}
\newcommand{\Expo}{\mathcal{E}}
\newcommand{\Poi}{\mathcal{P}}
\newcommand{\Ff}{\mathbb{F}}
\newcommand{\1}{\mathbf{1}}
\newcommand{\dd}{\,\mathrm{d}}
\DeclareMathOperator{\Var}{\mathbb{V}}
\DeclareMathOperator{\Cov}{Cov}
\DeclareMathOperator{\card}{card}
\DeclareMathOperator{\diag}{diag}
\newcommand{\indep}{\perp\!\!\!\perp}
\newcommand{\ps}{\ \PP\text{-p.s.}}
\newcommand{\tendloi}{\xrightarrow[n\to+\infty]{\text{loi}}}
\newcommand{\tendps}{\xrightarrow[n\to+\infty]{\text{p.s.}}}
\newcommand{\tendproba}{\xrightarrow[n\to+\infty]{\PP}}
\newcommand{\tendLp}{\xrightarrow[n\to+\infty]{L^p}}
\newcommand{\tend}{\xrightarrow[n\to+\infty]{}}

%     \fiche{C5}{Vecteurs gaussiens}{Gaussian vectors}
% ==============================================================
\usepackage[T1]{fontenc}
\usepackage[utf8]{inputenc}
\usepackage[french]{babel}
\usepackage{lmodern}
\usepackage[a4paper,margin=1.9cm,top=1.7cm,bottom=1.9cm]{geometry}
\usepackage{amsmath,amssymb,amsthm,mathtools}
\usepackage[inline]{enumitem}
\usepackage{booktabs,array,tabularx}
\usepackage{xcolor}
\usepackage[most]{tcolorbox}
\usepackage{fancyhdr}
\usepackage{hyperref}
\hypersetup{colorlinks=true,linkcolor=blue!50!black,urlcolor=blue!50!black}

\setlist{itemsep=1pt,topsep=3pt,parsep=0pt}
\setlength{\parindent}{0pt}
\setlength{\parskip}{3pt}
\emergencystretch=3em

% ---------- Identification de la fiche -------------------------
\newcommand{\ficheid}{}
\newcommand{\fichetitre}{}
\newcommand{\fiche}[3]{%
  \renewcommand{\ficheid}{#1}\renewcommand{\fichetitre}{#2}%
  \pagestyle{fancy}\fancyhf{}%
  \fancyhead[L]{\small\textsf{Advanced Probabilities -- Fiche #1}}%
  \fancyhead[R]{\small\textsf{#2}}%
  \fancyfoot[C]{\small\thepage}%
  \begin{center}
    {\LARGE\bfseries Fiche #1 -- #2}\par\smallskip
    {\large\itshape (#3)}\par\smallskip
    {\small Cours de V.\,M.~Lambert, M1 -- semestre 1}
  \end{center}\vspace{-2pt}\hrule\vspace{6pt}}

% ---------- Environnements numérotés (un seul compteur) --------
\theoremstyle{definition}
\newtheorem{definition}{Définition}
\renewcommand{\thedefinition}{\ficheid.\arabic{definition}}
\newtheorem{theoreme}[definition]{Théorème}
\newtheorem{proposition}[definition]{Proposition}
\newtheorem{lemme}[definition]{Lemme}
\newtheorem{corollaire}[definition]{Corollaire}
\newtheorem{exemple}[definition]{Exemple}
\theoremstyle{remark}
\newtheorem{remarque}[definition]{Remarque}

% Résultat hors cours (standard, signalé comme tel)
\theoremstyle{definition}
\newtheorem{horscours}[definition]{Hors cours}

% Idée de preuve
\newenvironment{idee}{\par\small\textit{Idée de preuve.}~}{\par}

% ---------- Encadrés -----------------------------------------
\tcbset{fichebox/.style={enhanced,breakable,boxrule=0.6pt,arc=2pt,
  left=5pt,right=5pt,top=3pt,bottom=3pt,fonttitle=\bfseries}}
\newtcolorbox{piege}[1][]{fichebox,colback=red!4,colframe=red!60!black,
  title={Pièges classiques},#1}
\newtcolorbox{liens}[1][]{fichebox,colback=blue!3,colframe=blue!50!black,
  title={Liens entre notions},#1}
\newtcolorbox{methode}[1][Méthode]{fichebox,colback=green!4,colframe=green!40!black,title={#1}}
\newtcolorbox{alerte}[1][]{fichebox,colback=orange!6,colframe=orange!70!black,
  title={Attention à l'énoncé},#1}

% Exercice de TD (fiches de rappels)
\newcounter{exo}
\newenvironment{exo}[1][]{\refstepcounter{exo}\par\medskip
  \noindent{\large\bfseries Exercice~\theexo}\ifx\relax#1\relax\else~{\large\bfseries -- #1}\fi\par\nopagebreak\smallskip}{\par}
\newcommand{\outils}{\textbf{Outils du cours.}~}
\newcommand{\meth}{\textbf{Méthode / astuce.}~}
\newcommand{\indic}{\textbf{Indication.}~}
\newcommand{\rep}{\textbf{Réponse.}~}

% ---------- Notations du cours ---------------------------------
\newcommand{\R}{\mathbb{R}}
\newcommand{\N}{\mathbb{N}}
\newcommand{\Q}{\mathbb{Q}}
\newcommand{\C}{\mathbb{C}}
\newcommand{\PP}{\mathbb{P}}
\newcommand{\QQ}{\mathbb{Q}}
\newcommand{\E}{\mathbb{E}}
\newcommand{\F}{\mathcal{F}}
\newcommand{\G}{\mathcal{G}}
\newcommand{\Hc}{\mathcal{H}}
\newcommand{\Tc}{\mathcal{T}}
\newcommand{\Bc}{\mathcal{B}}
\newcommand{\Cc}{\mathcal{C}}
\newcommand{\Ec}{\mathcal{E}}
\newcommand{\Pc}{\mathcal{P}}
\newcommand{\Nc}{\mathcal{N}}
\newcommand{\LN}{\mathcal{LN}}
\newcommand{\Expo}{\mathcal{E}}
\newcommand{\Poi}{\mathcal{P}}
\newcommand{\Ff}{\mathbb{F}}
\newcommand{\1}{\mathbf{1}}
\newcommand{\dd}{\,\mathrm{d}}
\DeclareMathOperator{\Var}{\mathbb{V}}
\DeclareMathOperator{\Cov}{Cov}
\DeclareMathOperator{\card}{card}
\DeclareMathOperator{\diag}{diag}
\newcommand{\indep}{\perp\!\!\!\perp}
\newcommand{\ps}{\ \PP\text{-p.s.}}
\newcommand{\tendloi}{\xrightarrow[n\to+\infty]{\text{loi}}}
\newcommand{\tendps}{\xrightarrow[n\to+\infty]{\text{p.s.}}}
\newcommand{\tendproba}{\xrightarrow[n\to+\infty]{\PP}}
\newcommand{\tendLp}{\xrightarrow[n\to+\infty]{L^p}}
\newcommand{\tend}{\xrightarrow[n\to+\infty]{}}

\begin{document}
\fiche{C6}{Espérance conditionnelle}{Conditional expectation}

\textit{Sources : slides, chapitre 7 (Déf.~50--52, Thm~12--15) ; poly « Conditional Expectation » (Déf.~1.1, Prop.~2.1--2.7,
Lemme~2.1, Thm~2.1) ; poly « Gaussian vector » (Déf.~14--17, Thm~18--21). Cadre : $(\Omega,\F,\PP)$, $\G,\Hc$ sous-tribus de $\F$.}

\section*{1. Conditionnement élémentaire}

\begin{definition}[sachant un événement]\label{C6:evt}
Pour $X$ intégrable et $B\in\F$ avec $\PP(B)>0$ : $\E(X\mid B)=\dfrac{\E(X\1_B)}{\PP(B)}$ (un \emph{nombre}).
Avec $X=\1_A$ : $\E(\1_A\mid B)=\PP(A\mid B)$.
\end{definition}

\begin{definition}[sachant une v.a.\ discrète]\label{C6:discret}
Soit $X$ réelle intégrable et $Y$ prenant un nombre au plus dénombrable de valeurs. $\E(X\mid Y)$ est la \emph{v.a.}
$\E(X\mid Y)(\omega)=\phi(Y(\omega))$, où
$\phi(y)=\E(X\mid Y=y)=\dfrac{\E(X\1_{\{Y=y\}})}{\PP(Y=y)}$ si $\PP(Y=y)>0$, et $\phi(y)=0$ sinon.
\end{definition}

\section*{2. Définition générale}

\begin{definition}[Espérance conditionnelle sachant une tribu]\label{C6:def}
Soit $Y\in L^1(\Omega,\F,\PP)$ et $\G\subset\F$ une sous-tribu. L'espérance conditionnelle de $Y$ sachant $\G$ est \emph{toute}
v.a.\ $Z$ telle que
\begin{enumerate}[label=(\roman*)]
  \item $Z$ est \textbf{$\G$-mesurable} ;
  \item $\forall G\in\G$, $\E(\1_GY)=\E(\1_GZ)$ \quad (i.e.\ $\int_GY\dd\PP=\int_GZ\dd\PP$).
\end{enumerate}
Une telle $Z$ existe et est \textbf{unique à égalité $\PP$-p.s.\ près} ; on note $Z=\E(Y\mid\G)$.
\end{definition}

\begin{remarque}[compléments du cours]\leavevmode
\begin{itemize}
  \item (ii) équivaut à : $\E(XY)=\E(XZ)$ pour toute v.a.\ $X$ $\G$-mesurable (\emph{bornée}, pour que tout soit intégrable).
  \item Si $\Omega$ est fini et $\PP(\{\omega\})>0$ pour tout $\omega$, $\E(Y\mid\G)$ est unique (pas seulement p.s.).
  \item Existence : si $Y\in L^2$, $\E(Y\mid\G)$ est la \emph{projection orthogonale} de $Y$ sur le sous-espace $L^2(\Omega,\G,\PP)$,
  c'est-à-dire la meilleure approximation de $Y$ au sens quadratique par une v.a.\ $\G$-mesurable de carré intégrable ; cas $L^1$ par densité de $L^2$
  ou par Radon--Nikodym (fiche C7).
  \item Vecteur : $\E(Y\mid\G)=(\E(Y_1\mid\G),\dots,\E(Y_d\mid\G))$. Pour $X:\Omega\to(E,\Ec)$ : $\E(Y\mid X):=\E(Y\mid\sigma(X))$.
  \item Toute v.a.\ $\sigma(X)$-mesurable s'écrit $\varphi(X)$ avec $\varphi$ mesurable ; donc $\E(Y\mid X)=\varphi(X)$, et on note
  $\varphi(x)=\E(Y\mid X=x)$.
\end{itemize}
\end{remarque}

\section*{3. Propriétés (toutes vraies $\PP$-p.s.)}

\begin{proposition}\label{C6:prop}
Soient $X,Y\in L^1$ et $\alpha,\beta\in\R$.
\begin{enumerate}
  \item \textbf{Espérance totale} : $\E\bigl(\E(Y\mid\G)\bigr)=\E(Y)$.
  \item \textbf{Linéarité} : $\E(\alpha X+\beta Y\mid\G)=\alpha\E(X\mid\G)+\beta\E(Y\mid\G)$.
  \item \textbf{Positivité / monotonie} : $Y\ge0$ p.s.\ $\Rightarrow\E(Y\mid\G)\ge0$ ; $X\le Y$ p.s.\ $\Rightarrow\E(X\mid\G)\le\E(Y\mid\G)$.
  \item \textbf{Ce qui est connu} : si $Y$ est $\G$-mesurable, $\E(Y\mid\G)=Y$. Si $X$ est $\G$-mesurable et $\E|XY|<\infty$,
  $\E(XY\mid\G)=X\,\E(Y\mid\G)$.
  \item Si $X$ est $\G$-mesurable et $\E|XY|<\infty$ (alors $X\,\E(Y\mid\G)=\E(XY\mid\G)\in L^1$) : $\E\bigl(X\,\E(Y\mid\G)\bigr)=\E(XY)$.
  % À VÉRIFIER : le cours suppose seulement E|X E(Y|G)| < infini, ce qui ne garantit pas que E(XY) existe
  % (ex. Z dans L^1 hors L^2, eps = +-1 independant, Y = eps Z, G = sigma(Z), X = Z).
  (Le cours écrit l'hypothèse $\E|X\,\E(Y\mid\G)|<\infty$, insuffisante pour que $\E(XY)$ existe.)
  \item \textbf{Tour} (\emph{tower law, chaining rule}) : si $\Hc\subset\G$, $\E\bigl(\E(Y\mid\G)\mid\Hc\bigr)=\E(Y\mid\Hc)$.
  « La plus petite information l'emporte. »
  \item \textbf{Indépendance} : si $\sigma(Y)$ et $\G$ sont indépendantes, $\E(Y\mid\G)=\E(Y)$. Plus généralement, si $\G$ est
  indépendante de $\sigma(\sigma(Y)\cup\Hc)$, alors $\E(Y\mid\sigma(\Hc\cup\G))=\E(Y\mid\Hc)$.
  \item \textbf{Jensen conditionnel} : $\Phi$ convexe et $\E|\Phi(Y)|<\infty$ $\Rightarrow$ $\Phi(\E(Y\mid\G))\le\E(\Phi(Y)\mid\G)$.
\end{enumerate}
\end{proposition}
\begin{idee}
On vérifie (i)--(ii) pour le candidat puis on invoque l'unicité p.s.
1 : prendre $G=\Omega$. 2 : $\alpha\E(X\mid\G)+\beta\E(Y\mid\G)$ est $\G$-mesurable et
$\E\bigl((\alpha\E(X\mid\G)+\beta\E(Y\mid\G))\1_G\bigr)=\alpha\E(X\1_G)+\beta\E(Y\1_G)$.
6 : $\E(Y\mid\Hc)$ est $\Hc$-mesurable et pour $H\in\Hc\subset\G$, $\E(\1_H\E(Y\mid\G))=\E(\1_HY)=\E(\1_H\E(Y\mid\Hc))$.
7 : $\E(\1_GY)=\PP(G)\E(Y)=\E(\1_G\E(Y))$ par indépendance.
\end{idee}

\begin{proposition}[gel / \emph{freezing}, slides]\label{C6:gel}
Soit $X:(\Omega,\F)\to(E,\Ec)$ \textbf{indépendante} de $\G$, $Y$ \textbf{$\G$-mesurable} réelle, et $\Phi:\R\times E\to\R$
$\Bc(\R)\otimes\Ec$-mesurable, positive ou bornée. Alors $\varphi(y):=\E\,\Phi(y,X)$ est mesurable et
$\E(\Phi(Y,X)\mid\G)=\varphi(Y)$ p.s. (« on gèle $Y$ comme une constante, on intègre en $X$ »).
\end{proposition}

\begin{proposition}[convergences conditionnelles, poly]\label{C6:conv}\leavevmode
\begin{enumerate}
  \item (\emph{convergence monotone}) $0\le X_n\uparrow X$ p.s.\ $\Rightarrow\E(X_n\mid\G)\uparrow\E(X\mid\G)$ p.s.
  \item (\emph{Fatou}) $X_n\ge0$ $\Rightarrow\E(\liminf X_n\mid\G)\le\liminf\E(X_n\mid\G)$ p.s.
  \item (\emph{convergence dominée}) $X_n\to X$ p.s.\ et $|X_n|\le Y\in L^1$ $\Rightarrow\E(X_n\mid\G)\to\E(X\mid\G)$ p.s.
\end{enumerate}
\end{proposition}

\section*{4. Calculs explicites}

\begin{theoreme}[partition finie]\label{C6:partition}
Si $X\in L^1$ et $\G=\sigma(A_1,\dots,A_n)$ avec $(A_j)$ partition de $\Omega$ (et $\PP(A_j)>0$), alors
$\displaystyle\E(X\mid\G)=\sum_{j=1}^n\frac{\E(X\1_{A_j})}{\PP(A_j)}\1_{A_j}$ p.s.
\end{theoreme}
\begin{idee}
$\G$-mesurable $\Rightarrow$ constante sur chaque $A_j$ : $\E(X\mid\G)=\sum_j\lambda_j\1_{A_j}$. Puis
$\E(\1_{A_i}X)=\E(\1_{A_i}\E(X\mid\G))=\lambda_i\PP(A_i)$.
\end{idee}

\begin{theoreme}[cas à densité]\label{C6:densite}
Soient $X$ et $Y$ des vecteurs aléatoires à valeurs dans $\R^p$ et $\R^k$ tels que $(X,Y)$ ait une densité $f(x,y)$ par rapport à Lebesgue sur $\R^{p+k}$, et
$f_X(x)=\int_{\R^k}f(x,y)\dd y$. La loi de $Y$ sachant $\{X=x\}$ a pour densité
\[f_Y^{X=x}(y)=\frac{f(x,y)}{f_X(x)}\quad\text{si }f_X(x)>0\qquad(\text{n'importe quelle densité si }f_X(x)=0).\]
Pour $g:\R^{p+k}\to\R$ mesurable, positive ou bornée : $\E(g(X,Y)\mid X=x)=\int_{\R^k}g(x,y)f_Y^{X=x}(y)\dd y=:\psi(x)$ et
$\E(g(X,Y)\mid X)=\psi(X)$ p.s.
\end{theoreme}

\begin{theoreme}[cas gaussien : régression linéaire]\label{C6:gauss}
Si $(Y_1,\dots,Y_n,X)$ est un \textbf{vecteur gaussien} (et $\Sigma_Y$ inversible), alors
\begin{gather*}\E(X\mid\sigma(Y_1,\dots,Y_n))=\E(X)+\sum_{j=1}^n\lambda_j(Y_j-\E Y_j)=\E(X)+\Lambda^t(Y-\E Y)\quad\text{p.s.},\\
\Lambda=\Sigma_Y^{-1}\begin{pmatrix}\Cov(X,Y_1)\\\vdots\\\Cov(X,Y_n)\end{pmatrix}.\end{gather*}
L'espérance conditionnelle coïncide avec la régression linéaire de $X$ sur $(Y_1,\dots,Y_n)$.
\end{theoreme}
\begin{idee}
Avec $Z=\E(X\mid Y)$ et $Y_i-\E Y_i$ $\sigma(Y)$-mesurable : $\E(Z(Y_i-\E Y_i))=\E(X(Y_i-\E Y_i))=\Cov(X,Y_i)$ ; en injectant la forme
affine, $\sum_j\lambda_j\Cov(Y_j,Y_i)=\Cov(X,Y_i)$, i.e.\ $\Sigma_Y\Lambda=(\Cov(X,Y_i))_i$. La forme affine est admise dans le cours ; justification : avec ce $\Lambda$, $R=X-\E X-\Lambda^t(Y-\E Y)$ forme avec $Y$ un vecteur gaussien et $\Cov(R,Y_i)=0$, donc $R\indep Y$ (C5) et $\E(R\mid Y)=\E R=0$.
\end{idee}

\begin{horscours}[loi conditionnelle gaussienne]\label{C6:gauss-loi}
Sous les hypothèses du Thm~\ref{C6:gauss}, $R:=X-\E(X\mid Y)$ est indépendant de $Y$ (le vecteur $(R,Y)$ est gaussien et
$\Cov(R,Y_i)=0$, C5), et $R\sim\Nc(0,s^2)$ avec $s^2=\Var X-C^t\Sigma_Y^{-1}C$, $C=(\Cov(X,Y_i))_i$.
Par le lemme de gel : la loi de $X$ sachant $Y$ est $\Nc\bigl(\E(X\mid Y),\,s^2\bigr)$, et $\E\bigl((X-\E(X\mid Y))^2\mid Y\bigr)=s^2$.
\end{horscours}

\begin{liens}
\begin{itemize}
  \item $\E(Y\mid\{\emptyset,\Omega\})=\E(Y)$ ; $\E(Y\mid\F)=Y$ ; $\G$ « entre les deux » donne une v.a.\ intermédiaire.
  \item Déf.~\ref{C6:discret} et Thm~\ref{C6:partition} sont des cas particuliers de la Déf.~\ref{C6:def} (avec $\G=\sigma(Y)$,
  tribu d'une partition).
  \item $\sigma(Y)\indep\G\Rightarrow\E(Y\mid\G)=\E Y$ ; réciproque fausse : $Y\sim\Nc(0,1)$, $\G=\sigma(Y^2)$, alors
  $\E(Y\mid\G)=0=\E Y$ ($(Y,Y^2)$ et $(-Y,Y^2)$ ont même loi) mais $Y$ n'est pas indépendante de $Y^2$.
  \item Pour $X,Y\in L^2$ : indépendance $\Rightarrow\E(X\mid Y)=\E X\Rightarrow\Cov(X,Y)=0$ ; aucune réciproque en général. Par exemple $Y\sim\Nc(0,1)$, $X=Y^2$ : $\Cov(X,Y)=\E Y^3=0$ mais $\E(X\mid Y)=Y^2\neq1=\E X$.
\end{itemize}
\end{liens}

\begin{piege}
\begin{itemize}
  \item $\E(Y\mid\G)$ est une \textbf{variable aléatoire} ($\G$-mesurable), pas un nombre ; toutes les égalités sont \textbf{p.s.}
  \item Hypothèse d'\textbf{intégrabilité} : $Y\in L^1$ (et $XY\in L^1$ pour sortir $X$).
  \item On ne sort que ce qui est \textbf{$\G$-mesurable} : $\E(XY\mid\G)\neq X\E(Y\mid\G)$ en général.
  \item Tour : $\Hc\subset\G$ ; dans les deux ordres on obtient $\E(Y\mid\Hc)$ (la plus petite tribu gagne) ; dans l'ordre $\E(\E(Y\mid\Hc)\mid\G)$, c'est « ce qui est connu », car $\E(Y\mid\Hc)$ est $\G$-mesurable.
  \item Jensen : sens $\Phi(\E(Y\mid\G))\le\E(\Phi(Y)\mid\G)$ pour $\Phi$ \emph{convexe} ; $\E(\Phi(Y)\mid\G)\neq\Phi(\E(Y\mid\G))$.
  \item Conditionner par $\{X=x\}$ quand $\PP(X=x)=0$ : passer par les densités (Thm~\ref{C6:densite}), pas par la Déf.~\ref{C6:evt}.
  \item Lemme de gel : exige $X\indep\G$ \emph{et} $Y$ $\G$-mesurable.
  \item Cas gaussien : la formule linéaire n'est valable que si le vecteur $(Y,X)$ est \emph{gaussien}.
\end{itemize}
\end{piege}

\section*{Exemples types}

\begin{exemple}[dé, slides]\label{C6:ex-de}
$\Omega=\{1,\dots,6\}$ uniforme, $X(\omega)=\omega$, $Y=\1_{\{\omega\text{ impair}\}}$. Sur $\{Y=1\}$ :
$\E(X\mid Y)=\frac{\frac16(1+3+5)}{3/6}=3$ ; sur $\{Y=0\}$ : $\frac{\frac16(2+4+6)}{3/6}=4$. Donc $\E(X\mid Y)=4-Y$.
Avec $Z=1$ sur $\{1,3\}$, $2$ en $5$, $3$ sur $\{2,4\}$, $4$ en $6$ (plus informative : $\sigma(Y)\subset\sigma(Z)$), la tour donne $\E(\E(X\mid Z)\mid Y)=\E(X\mid Y)$ (vérifié dans les slides).
\end{exemple}

\begin{exemple}[arbre binomial, slides]\label{C6:ex-arbre}
Modèle C1 ($0<p<1$), $S_2(UU)=s_0u^2$, $S_2(UD)=S_2(DU)=s_0ud$, $S_2(DD)=s_0d^2$, $\F_1=\sigma(A_U)$. Par le Thm~\ref{C6:partition} :
$\E(S_2\1_{A_U})=p^2s_0u^2+p(1-p)s_0ud=ps_0u(pu+(1-p)d)$ et $\PP(A_U)=p$, idem sur $A_D$ ; d'où
$\E(S_2\mid\F_1)=(pu+(1-p)d)(s_0u\1_{A_U}+s_0d\1_{A_D})=(pu+(1-p)d)\,S_1$ p.s.
\end{exemple}

\begin{exemple}[densité]\label{C6:ex-densite}
$(X,Y)$ de densité $f(x,y)=e^{-y}\1_{0<x<y}$. $f_X(x)=\int_x^{\infty}e^{-y}\dd y=e^{-x}\1_{x>0}$ et
$f_Y^{X=x}(y)=e^{-(y-x)}\1_{y>x}$ : sachant $X=x$, $Y-x\sim\Expo(1)$. Donc $\E(Y\mid X=x)=x+1$ et $\E(Y\mid X)=X+1$ p.s.
\end{exemple}

\end{document}
